% Gauss, Werke X/1 (1917), native87–89, printed80–82.
\SourceLeaf{087}\NoSourceHead
\SourceGroup{II}
\HistoricTitle{[ÜBER DIE ANZAHL DER ZERLEGUNGEN EINER PRIMZAHL IN DREI QUADRATE.]}
\HistoricCitation{[Ein Zettel in Ea 5, Kapsel 40.]}
\par\indent Wenn $A$ eine Primzahl ist, so gibt es so viele Formen der Divisoren
$4n+1$ von $\square+A\square$ als es verschiedene Arten gibt $A$ in drei Quadrate zu
zerlegen.
\par\medskip\noindent\hspace*{\parindent}\textls[100]{Bew[eis.]}
\par\noindent\hspace*{\parindent}1.\quad\textit{Eine jede Zerlegung von $A$ in drei Quadrate gibt eine Form oder einen Multiplicator.}

Ist nemlich
\[
A=aa+bb+cc,
\]
\noindent so erhält man dadurch den Multiplicator $aa+bb$ bei $\dfrac{ca}{b}$ oder $\dfrac{cb}{a}$. —
\[
aa+bb;\ aa+cc;\ bb+cc
\]
\noindent sind verbunden und geben folglich nur eine Form.
\par\noindent\hspace*{\parindent}2.\quad\textit{Jeder Multiplicator gibt eine Zerlegung von $A$.}

Man darf voraussetzen, dass das von allen Zahlen kleiner als $A$ gilt
(vermöge der Depressionsmethode), dass also der Multiplicator $(=\triangle)=mm+nn+pp$, so dass $mm+nn$ ein Multiplicator von $A$ oder
\[
(mm+nn)A=xx+(mm+nn+pp)yy
\]
\noindent also
\[
A=yy+\frac{ppyy+xx}{mm+nn}=yy+zz+ww.
\]
\clearpage
\SourceLeaf{088}\ArithmeticHead{81}{R}{ZERLEGUNGEN EINER PRIMZAHL IN DREI QUADRATE.}
\noindent Beide Sätze gelten nicht bloss von Primzahlen, sondern auch von zusammengesetzten. (Zus. $(zz+ww)(mm+nn+pp)=xx+App.$)

3.\quad Es bleibt uns also zu beweisen übrig, dass wenn $A$ eine Primzahl
ist, kein Multiplicator (im eingeschränkten Sinne) zu mehr als einer Zerlegung
gehören könne. Bei zusammengesetzten Zahlen ist dies nicht immer der Fall,
z. B. die Zerlegungen von 185 in
\[
4+81+100\ \text{und}\ 36+49+100
\]
\noindent geben einen Multiplicator 85 bei $\dfrac{10\mathpunct{.}9}{2}$ und $\dfrac{10\mathpunct{.}7}{6}$, beide gleich 45. Der kleinste
Multiplicator ist 9, so dass
\[
1\mathpunct{.}2+2\mathpunct{.}9-2\mathpunct{.}10=0\ \text{und}\ 1\mathpunct{.}6+2\mathpunct{.}7-2\mathpunct{.}10=0.
\]
\par\noindent\hspace*{\parindent}4.\quad Es sei also
\[
A=aa+bb+cc=a'a'+b'b'+c'c'
\]
\noindent und beide Zerlegungen sollen zusammenhangende Multiplicatoren geben, also
\[
(aa+bb)(a'a'+b'b')=\square+A\square
\]
\noindent oder wenn
\[
\begin{gathered}
a'a'+b'b'=mm+nn+pp=\triangle,\\
am+bn+cp=0
\end{gathered}
\]
\noindent (weil
\[
a'a'+b'b'=(aa+bb)xx+2qxy+ryy=\square+\square+\square).
\]
\par\noindent\hspace*{\parindent}Daher müsste sein
\[
(mm+nn)A=\square+\triangle\square,
\]
\noindent und
\[
mm+nn=\square+\triangle\square
\]
\noindent (weil $A=\square+\triangle\square$ eine Primzahl), welches nur möglich ist, wenn
\[
mm+nn=\square;
\]
\noindent wir müssen also gegenwärtig setzen
\[
a'a'+b'b'=ff+gg,
\]
\par\vfill\noindent X$_1$.\hfill11\par\clearpage
\SourceLeaf{089}\ArithmeticHead{82}{L}{ARITHMETIK.\quad NACHLASS.}
\noindent so dass
\[
af+bg=0,
\]
\noindent also wenn
\[
\begin{gathered}
\frac ab=\frac{\mu t}{\mu u},\\
f=\nu u;\ g=\nu t,
\end{gathered}
\]
\noindent und so hätten wir endlich für $A$ die beiden Zerlegungen
\[
\mu\mu tt+\mu\mu uu+cc
\]
\noindent und
\[
\nu\nu tt+\nu\nu uu+c'c',
\]
\noindent welches nicht sein kann, wenn $A$ eine Primzahl ist; daher
\[
\mu=\nu
\]
\noindent und beide Zerlegungen dieselben sind.
\begin{center}\textls[100]{Erläuterungen.}\end{center}
\par\indent 1) Wenn $aa+bb$ und $a'a'+b'b'\ (=\triangle)$ zusammenhangen, so ist
\[
(aa+bb)\triangle=xx+(aa+bb+cc)yy,
\]
\noindent woraus sogleich eine Zerlegung von $\triangle$ in drei Quadrate folgt, wovon das
eine $=yy$. Man wird leicht finden, dass diese Quadrate
\[
mm+nn+pp
\]
\noindent von der Art sind, dass
\[
am+bn+cp=0.
\]
\par\noindent\hspace*{\parindent}2) Wenn
\[
aa'+bb'+cc'=0,
\]
\noindent so kann man allemal setzen
\[
(aa+bb+cc)(a'a'+b'b')=AA+(a'a'+b'b'+c'c')BB
\]
\noindent und
\[
(a'a'+b'b'+c'c')(aa+bb)=A'A'+(aa+bb+cc)B'B';
\]
\noindent man setze
\[
B=c;\ B'=c'\ \text{und}\ A=ab'-a'b=A'.
\]
\ClosingRule\clearpage
